% ============================================================
% Section 90: 钾金属低温电子热容与费米参数
% ============================================================
\section{固体理论：钾金属低温电子热容与费米参数}
\label{sec:potassium-specific-heat}

\begin{modelbox}[题目：由低温电子热容实验估算钾的费米参数]
低温下金属钾的摩尔电子热容实验结果为 $c_e=2.08T\,\text{mJ/(mol}\cdot\text{K)}$。X射线衍射测得晶格常数 $a=5.333$ \AA。在自由电子气模型下估算：
\begin{enumerate}
    \item 费米温度 $T_F$；
    \item 导带电子有效质量 $m^*$；
    \item 费米速度 $v_F$；
    \item 费米面上的能态密度 $\rho(E_F)$。
\end{enumerate}
\end{modelbox}

\begin{tipbox}[考点快速分析]
\begin{itemize}
    \item \textbf{低温电子比热}：$c_e=\dfrac{\pi^2}{2}ZR\dfrac{T}{T_F}$，线性温度依赖
    \item \textbf{费米温度}：由比热系数直接提取
    \item \textbf{有效质量}：$m^*=\hbar^2k_F^2/(2k_BT_F)$，反映能带修正
    \item \textbf{态密度}：$\rho(E_F)=3n/(2E_F)$，决定输运和热力学性质
\end{itemize}
\end{tipbox}

\begin{derivationbox}[标准解题过程]
\textbf{Step 1: 费米温度}

自由电子气模型中，低温摩尔电子热容
\begin{equation}
c_e=\frac{\pi^2}{2}ZR\frac{T}{T_F},
\end{equation}
其中 $Z=1$（钾为单价金属），$R=8.314\,\text{J/(mol}\cdot\text{K)}$。

已知 $c_e/T=2.08\times10^{-3}\,\text{J/(mol}\cdot\text{K}^2)$，故
\begin{equation}
\boxed{T_F=\frac{\pi^2R}{2\times2.08\times10^{-3}}
=\frac{9.87\times8.314}{4.16\times10^{-3}}
\approx1.97\times10^4\,\text{K}.}
\end{equation}

\textbf{Step 2: 导带电子有效质量}

钾为 BCC 结构，$n=2/a^3$。费米波矢
\begin{equation}
k_F=(3\pi^2n)^{1/3}=\frac{(6\pi^2)^{1/3}}{a}
\approx\frac{3.898}{5.333\times10^{-10}}
\approx7.31\times10^9\,\text{m}^{-1}.
\end{equation}

由 $E_F=\hbar^2k_F^2/(2m^*)=k_BT_F$，解出
\begin{equation}
\boxed{m^*=\frac{\hbar^2k_F^2}{2k_BT_F}
\approx\frac{(1.055\times10^{-34})^2\times(7.31\times10^9)^2}{2\times1.38\times10^{-23}\times1.97\times10^4}
\approx1.09\times10^{-30}\,\text{kg}\approx1.2\,m_0.}
\end{equation}

\textbf{Step 3: 费米速度}

\begin{equation}
\boxed{v_F=\frac{\hbar k_F}{m^*}
\approx\frac{1.055\times10^{-34}\times7.31\times10^9}{1.09\times10^{-30}}
\approx7.05\times10^5\,\text{m/s}=0.705\times10^8\,\text{cm/s}.}
\end{equation}

\textbf{Step 4: 费米面态密度}

三维自由电子气态密度
\begin{equation}
\rho(E)=\frac{1}{2\pi^2}\Bigl(\frac{2m^*}{\hbar^2}\Bigr)^{3/2}E^{1/2}.
\end{equation}

在费米面处
\begin{align}
\rho(E_F)&=\frac{1}{2\pi^2}\Bigl(\frac{2m^*}{\hbar^2}\Bigr)^{3/2}(k_BT_F)^{1/2} \notag\\
&\approx7.28\times10^{46}\,\text{J}^{-1}\text{m}^{-3} \notag\\
&\approx\boxed{1.17\times10^{28}\,\text{eV}^{-1}\text{m}^{-3}.}
\end{align}
\end{derivationbox}

\begin{formulabox}[答案汇总]
\begin{center}
\begin{tabular}{ll}
\toprule
\textbf{物理量} & \textbf{数值} \\
\midrule
费米温度 $T_F$ & $1.97\times10^4\,\text{K}$ \\
有效质量 $m^*$ & $1.2\,m_0$ \\
费米速度 $v_F$ & $0.705\times10^8\,\text{cm/s}$ \\
态密度 $\rho(E_F)$ & $1.17\times10^{28}\,\text{eV}^{-1}\text{m}^{-3}$ \\
\bottomrule
\end{tabular}
\end{center}
\end{formulabox}

\begin{insightbox}[物理图像]
\begin{itemize}
    \item 低温电子比热是测定金属费米能和有效质量的\textbf{标准实验手段}。钾的 $m^*\approx1.2m_0$ 接近自由电子质量，说明钾是近自由电子金属，电子--声子和电子--电子相互作用对能带的修正较小。
    \item 费米速度 $v_F\sim10^5\,$m/s 虽然远小于光速，但比室温下气体分子热运动速度快约两个数量级，体现了电子简并态的高动能特征。
    \item 费米面态密度 $\rho(E_F)$ 是决定金属电导率（$
ho(E_F)$ 越大，载流子越多）、热导率、泡利顺磁磁化率等输运性质的核心参数。
\end{itemize}
\end{insightbox}
